\documentclass[10pt,a4paper]{amsart}
\usepackage[margin=19mm]{geometry}
\usepackage{amsmath,amssymb,mathtools}
\usepackage[scr=rsfs,cal=euler]{mathalfa}
\usepackage{xfrac}
\usepackage[unicode,hidelinks,bookmarks=false]{hyperref}
\newcommand{\ie}{i.e.\ }
\newcommand{\cf}{cf.\ }
\newcommand{\resp}{respectively\ }
\newcommand{\Subsection}{\subsection}
\providecommand{\nobreakdash}{\nobreak\hbox{-}}
\setlength{\emergencystretch}{2em}
\multlinegap0pt
\usepackage{bm}
\usepackage{bbm}
\usepackage{dsfont}
\usepackage{xspace}
\usepackage{cancel}
\usepackage{mathtools}
\usepackage{amsthm}
\makeatletter
\@ifundefined{theorem}{\newtheorem{theorem}{Theorem}}{}
\@ifundefined{proposition}{\newtheorem{proposition}[theorem]{Proposition}}{}
\makeatother
\newcommand{\CM}{\mathrm{CM}}
\newcommand{\Ext}{\mathrm{Ext}}
\newcommand{\ann}{\mathrm{ann}}
\newcommand{\module}{\mathrm{mod}}
\newcommand{\sann}{\underline{\mathrm{ann}}}
\title[Annihilation of cohomology over one dimensional almost Goren]{Annihilation of cohomology over one dimensional almost Gorenstein rings}
\author{\"Ozg\"ur Esentepe}
\date{Source: arXiv:2510.05720v1 (CC BY 4.0)}
\begin{document}
\maketitle

\noindent\emph{Source and licence.} Annihilation of cohomology over one dimensional almost Gorenstein rings. Version arXiv:2510.05720v1 (CC BY 4.0), distributed under the Creative Commons Attribution 4.0 International licence, as recorded in the arXiv licence metadata for this exact version and shown on the abstract page https://arxiv.org/abs/2510.05720v1. Full rights notice: licenses/arxiv-2510.05720-CC-BY-4.0.txt.\par\smallskip

\noindent\textit{Editorial scope.} Counted unit: the Proposition whose source label is \texttt{None} in the article (source0.tex lines 432--438, source label \texttt{None}), delivered together with the article's own complete original proof of it (source0.tex lines 439--455, one \texttt{proof} environment of 17 lines). No mathematics is written, repaired or re-derived: statement and proof are the article's own text line for line. Required context retained uncounted: main-proposition (lines 355--360). Every definition, display and label of those lines is retained unchanged. Omitted from this package and not redistributed: 1--354, 361--431, 456--629. Nothing inside a retained excerpt is omitted or altered: every retained body is delivered exactly as the article's lines are written, so no omission receipt of any kind accompanies this package. Citations: the retained excerpts carry no citation command at all, so no bibliography entry of the article is transcribed and no reference is redistributed. Reference adaptation: none was needed and none was made; every reference inside the delivered unit resolves inside it, so no unresolved reference or citation is produced. Printed slips kept literal: no printed word, symbol, hypothesis or number of the counted statement or of its proof was altered. Layout and class: the article's own class is \texttt{amsart} with packages including inputenc, amsmath, amsthm, amsfonts, amssymb, tabu, latexsym, geometry, enumerate, hyperref, mathrsfs, tikz, parskip, tikz-cd; the wrapper is the lane's pinned \texttt{amsart} editorial wrapper at 10pt on A4, which restates the theorem-like environments the retained text needs and adds the article's own macro definitions that the retained text needs (\texttt{\textbackslash CM}, \texttt{\textbackslash Ext}, \texttt{\textbackslash ann}, \texttt{\textbackslash module}, \texttt{\textbackslash sann}), with extra wrapper packages (bm, bbm, dsfont, xspace, cancel, mathtools). The delivered numbering is the unit's own. Assets: no retained span contains a figure environment, a graphics-inclusion command, a PDF-page inclusion, a table or a TikZ picture; no image file is copied into this package, the package declares no asset dependency and no image file is redistributed with it. Licence: version arXiv:2510.05720v1 of this article is distributed under the Creative Commons Attribution 4.0 International licence, as recorded in the arXiv licence metadata for this exact version and shown on the abstract page https://arxiv.org/abs/2510.05720v1. A full rights notice is delivered at licenses/arxiv-2510.05720-CC-BY-4.0.txt. Characters: every retained excerpt is pure ASCII, so the delivered engine, XeLaTeX with its default Latin Modern text font, reports no missing character.
\begin{proposition}\label{main-proposition}
    Let $R$ be a Cohen-Macaulay local ring with canonical module $\omega$ and $M$ a maximal Cohen-Macaulay $R$-module. If $D\Omega M \in \Omega \CM(R)$, then there is an equality 
    \begin{align*}
    \sann_R(M) = \sann_R(\Omega M).
    \end{align*}
\end{proposition}

\begin{proposition}
    Let $R$ be Cohen-Macaulay local ring with canonical module $\omega$ and $M$ be a maximal Cohen-Macaulay $R$-module. If both $M$ and $DM$ belong to $\Omega \CM(R)$, then there exists an equality
    \begin{align*}
        \bigcap_{X \in \CM(R)} \ann_R \Ext_R^{>0}(M,X) = \bigcap_{X \in \CM(R)} \ann_R \Ext_R^{>0}(X,M)
    \end{align*}
    of ideals.
\end{proposition}

\begin{proof}
    We note that as $M \in \CM(R)$, we also have that $\Omega_RM \in \CM(R)$. Hence, there is an inclusion 
    \begin{align*}
    \bigcap_{X \in \CM(R)} \ann_R \Ext_R^{>0}(M,X) \subseteq \ann_R\Ext_R^1(M, \Omega_R M) = \sann_R(M).
    \end{align*}
    On the other hand as $\CM(R) \subseteq \module(R)$, we have another inclusion 
    \begin{align*}
    \sann_R(M) = \bigcap_{X \in \module(R)} \ann_R \Ext_R^{>0}(M,X) \subseteq \bigcap_{X \in \CM(R)} \ann_R \Ext_R^{>0}(M,X)
    \end{align*}
    which means that the intersection on the left hand side of the equality in the statement is just the stable annihilator of $M$. 
    
    Since $D$ is a duality on $\CM(R)$, we also have equalities
    \begin{align*}
        \bigcap_{X \in \CM(R)} \ann_R \Ext_R^{>0}(X,M) = \bigcap_{X \in \CM(R)} \ann_R \Ext_R^{>0}(DX,M) = \bigcap_{X \in \CM(R)} \ann_R \Ext_R^{>0}(DM,X).
    \end{align*}
    Hence, the right hand side of the equality in the statement is $\sann_R(DM)$. Now, we are done by the proof of Proposition \ref{main-proposition}.
\end{proof}

\end{document}
